\begin{lemma}[A super-sequence is not both perfect and bad]
No super-sequence is simultaneously perfect and bad for the same relation.
\end{lemma}
\begin{proof}
Assume, for a contradiction, that

\[
h_P:\mathsf{PerfectSuperSequence}(r,f)
\qquad\text{and}\qquad
h_B:\mathsf{BadSuperSequence}(r,f).
\]

The structure \noderef{SuperSequence} records the front
$f.\mathsf{front}:\mathsf{Front}(F,X)$.  Its infinite-base field therefore
gives

\[
h_X:X.\mathsf{Infinite}.
\]

Apply \noderef{InfiniteSetEnumeration} to $X$ and $h_X$.  We obtain an
increasing enumeration $x\in\mathsf{IncSeq}$ and an equality

\[
h_{\mathrm{range}}:\operatorname{range}(x)=X.
\]

Consequently

\[
h_{x,X}:\operatorname{range}(x)\subseteq X,
\]

because if $n\in\operatorname{range}(x)$, then
$n\in X$ after rewriting by $h_{\mathrm{range}}$.

Apply \noderef{FrontPrefixExists} to the front $f.\mathsf{front}$, the
enumeration $x$, and $h_{x,X}$.  Choose $s$ and

\[
p_s:\mathsf{FrontPrefix}(F,s,x).
\]

By the definition of \noderef{FrontPrefix}, its two components are

\[
h_s:p_s.1\;\text{with }h_s:s\in F,
\qquad
q_s:p_s.2\;\text{with }
\mathsf{ProperPrefixSet}(s,\operatorname{range}(x)).
\tag{1}
\]

We next obtain the corresponding prefix after shifting.  Form the base
enumeration

\[
\bar x:=\langle x,h_{x,X}\rangle\in\mathsf{BaseIncSeq}(X),
\]

where the subtype is the one defined in \noderef{BaseIncSeq}.  By
\noderef{BaseShift} and the composition identity in \noderef{ShiftMap},
$\mathsf{BaseShift}_X(\bar x)$ is again in
$\mathsf{BaseIncSeq}(X)$, its first component is
$\mathsf{ShiftMap}(x)$, and its second component gives

\[
h_{\operatorname{shift}}:\operatorname{range}(\mathsf{ShiftMap}(x))\subseteq X.
\tag{2}
\]

Apply \noderef{FrontPrefixExists} again, now to $f.\mathsf{front}$, the
increasing enumeration $\mathsf{ShiftMap}(x)$, and (2).  Choose $t$ and

\[
p_t:\mathsf{FrontPrefix}(F,t,\mathsf{ShiftMap}(x)).
\]

Unpacking \noderef{FrontPrefix} once more gives

\[
h_t:p_t.1\;\text{with }h_t:t\in F,
\qquad
q_t:p_t.2\;\text{with }
\mathsf{ProperPrefixSet}
   (t,\operatorname{range}(\mathsf{ShiftMap}(x))).
\tag{3}
\]

The same increasing enumeration $x$ witnesses the finite shift from $s$ to
$t$: the first required proper-prefix condition is $q_s$ from (1), and the
second is $q_t$ from (3).  Hence, by the defining existential in
\noderef{FiniteShift},

\[
h_{\mathrm{shift}}^{s,t}:\mathsf{FiniteShift}(s,t).
\tag{4}
\]

Now use $h_P$ with the explicit front-membership proofs $h_s$ and $h_t$
and the witness (4).  By the defining statement in
\noderef{PerfectSuperSequence}, this yields

\[
r\bigl(f.\mathsf{value}\langle s,h_s\rangle,
       f.\mathsf{value}\langle t,h_t\rangle\bigr).
\tag{5}
\]

Use $h_B$ with exactly the same $s,t,h_s,h_t$, and (4).  The defining
statement in \noderef{BadSuperSequence} yields

\[
\neg r\bigl(f.\mathsf{value}\langle s,h_s\rangle,
            f.\mathsf{value}\langle t,h_t\rangle\bigr).
\tag{6}
\]

Statements (5) and (6) are a proposition and its negation, so they give a
contradiction.  Therefore no $f$ can satisfy both perfectness and badness.
\end{proof}
